\subsection{A mixed product with a uniform exact rank formula}
\label{sec:DE-A3-mixed-generic}

The following statement uses one fixed perfect field and one fixed embedded
surface in every degree.  Specialization is used to prove the formula;
the ordinary elliptic factor is not chosen anew as the degree changes.

\begin{theorem}[The geometric generic ordinary elliptic factor]
\label{thm:DE-A3-mixed-generic}
\status{Proved; complete independent root and monodromy proof audits passed.}
Let $k=\mathbf F_2(t)^{\mathrm{perf}}$ be the perfect closure of the
rational function field, and define
\[
 E_t:\ y^2+txy+y=x^3,\qquad
 E_0:\ y^2+y=x^3,\qquad
 X=E_t\times_k E_0,
 \qquad L=\OO_{E_t}(3O)\boxtimes\OO_{E_0}(3O).
\]
Use the cubic embeddings and the Segre map to embed $X$ in
$\mathbf P^8_k$.  The first elliptic factor is ordinary and the second
is supersingular.  For every $d\ge1$, the geometric generic relative
stable rank is
\[
 \boxed{\quad
 \dim E(Y,\OO_Y)_{\ss,\mathrm{gen}}
  =\begin{cases}
       9d^2-1,& d\ \text{odd},\\
       9d^2-2,& d\ \text{even}.
    \end{cases}\quad}                                      \tag{\ddagger}
\]
The relative coherent dimension is always $9d^2-1$.
Thus its generic Frobenius is bijective exactly in odd degree.
The full generic curve defect $g(Y)-f_Y$ is one in odd degree and
two in even degree.

The same formulas hold after replacing $E_t$ by any elliptic curve
over a perfect characteristic-two field whose $j$-invariant is
transcendental over $\mathbf F_2$, and taking geometric generic ranks.
This last assertion does not include arbitrary algebraic
$j$-invariants.
\end{theorem}
\begin{proof}
Consider the smooth elliptic scheme
\[
 \mathcal E\longrightarrow
 S=\Spec\mathbf F_2[t,(t^3+1)^{-1}]
\]
given by the same Weierstrass equation.  The integral formulas with
$a_1=t,a_3=1$ and the other coefficients zero give
$c_4=t^4-24t$ and $\Delta=t^3-27$; in characteristic two these are
$t^4$ and $t^3+1$.  The Hasse coefficient is $a_1=t$.
Thus the generic fibre is ordinary, while the fibre at $t=0$ is
exactly $E_0$.  In particular this is a specialization within smooth
elliptic schemes.

Fix an odd $d\ge1$ and put $m=3d$.  On
$\mathcal X=\mathcal E\times_S(E_0\times S)$ use
$\mathcal N_m=\OO_{\mathcal E}(mO)\boxtimes\OO_{E_0}(mO)$.
Its direct image has rank $m^2$ and commutes with base change,
because both elliptic factor bundles have positive degree and no
higher cohomology.  Let $\mathcal P_m\to S$ be its projective space
of nonzero sections up to scale, and let $\mathcal D_m$ be the
universal divisor.  The relative coherent quotient is a vector
bundle $\mathcal Q_m$ of rank $m^2-1$ on all of $\mathcal P_m$.
Indeed fibrewise K\"unneth gives
\[
 H^1(\mathcal X_s,\mathcal N_{m,s}^{-1})=0,\qquad
 h^2(\mathcal X_s,\mathcal N_{m,s}^{-1})=m^2,
 \qquad h^2(\mathcal X_s,\OO)=1.
\]
Multiplication by a nonzero section maps the middle cohomology
surjectively to the last one, since its Serre-dual map sends $1$
to that nonzero section.  Its kernel is exactly the relative
quotient, including for singular divisors.  This also proves
the required base-change and local-freeness assertions.

Absolute Frobenius gives an $\OO_{\mathcal P_m}$-linear map
$F_{\mathcal P_m}^*\mathcal Q_m\to\mathcal Q_m$.
Its invertibility locus is open.  On the fibre $t=0$, the
supersingular-product theorem
\ref{thm:DE-A3-ss-product-rank} supplies a smooth section where it
is invertible, since $d$ is odd.  The determinant is therefore
a nonzero section of a line bundle on the integral scheme
$\mathcal P_m$.  It cannot vanish on its entire generic fibre:
otherwise it would vanish at the generic point of
$\mathcal P_m$, hence everywhere.  The geometric generic curve
over $\mathbf F_2(t)$ consequently has relative rank $m^2-1$.
Passing to its perfect closure preserves this rank, as does any
further algebraically closed extension.

This argument applies to every odd $d$ at the same generic point
of $S$.  It makes no assertion that an infinite intersection of
nonempty open subsets of $S$ contains a closed point.  The
elliptic curve over the fixed field $k$ is therefore the same
$E_t$ for all of these conclusions.  The smooth-member open is
nonempty in each degree because both factor line bundles have
degree at least three.

Now work over $\bar k$ and let $Y\in|L^d|$ be smooth, with $m=3d$.
The Picard embedding and connected-Prym proof of
Lemma~\ref{lem:DE-A3-ss-prym} applies word for word to
$E_t\times E_0$: it only used negative-degree K\"unneth vanishing
for the embedding, and the diagonal product polarization for
$qj=[m]$.  Thus
\[
 0\longrightarrow P\longrightarrow J_Y\xrightarrow{q}X
 \longrightarrow0
\]
is exact with $P$ an abelian variety.  Here $g(Y)=m^2+1$,
$f_X=1$, $f_Y=f_P+1$, and $\dim P=g(Y)-2$.
If $m$ is even, $j$ embeds $X[2]$ into $P[2]$.  The identity
component of $\ker V_X$ has order $2$, since $\dim X=2$ and
$f_X=1$.  Applying Frobenius twist and functorial finite-group
Verschiebung to that immersion embeds this connected group
into $\ker V_P$.  Therefore
\[
 \dim P-f_P\ge1,\qquad f_Y\le g(Y)-2.
\]
Exactness of stable parts, now subtracting the one-dimensional
ambient stable part, gives the upper bound
$\dim E(Y,\OO_Y)_{\ss}\le m^2-2$.

For even $d\ge2$, choose generic smooth members
$C\in|L|$ and $D\in|L^{d-1}|$ in the odd-degree maximal-rank
opens just proved.  Each has full genus-minus-$2$-rank defect
one, since its ambient stable rank is one.  They can be chosen
transverse, by very ampleness and Bertini.  The nodal union
$C\cup D\in|L^d|$ has full defect two: the normalization
sequence adds the graph rank bijectively for Frobenius and
adds the two component defects.  Constant coherent dimensions
and the open-rank argument in
Theorem~\ref{thm:DE-A3-ss-product-rank} deform it to a smooth
member of relative rank at least $m^2-2$.  The preceding upper
bound makes this equality.  This proves (\ddagger).
Surjectivity from ambient degree-$d$ hypersurfaces to
$H^0(X,L^d)$ follows from the two plane-cubic coordinate rings
and the Segre map exactly as in that theorem.

Finally the rational function
\[
                         j(E_t)=\frac{t^{12}}{t^3+1}
\]
is nonconstant.  If an elliptic curve $E$ has transcendental
$j(E)$, choose a root $u$ of
$u^{12}-j(E)(u^3+1)=0$ in an algebraically closed extension.
Then $u$ is transcendental over $\mathbf F_2$ and $u^3+1\ne0$.
The resulting inclusion $\mathbf F_2(t)\hookrightarrow\bar k$,
$t\mapsto u$, identifies $E$ geometrically with the base change
of $E_t$, since elliptic curves with the same $j$-invariant
are isomorphic over an algebraically closed field.  The origins
and their cubic line bundles are respected.  Geometric generic
ranks are invariant under this field extension, proving the
last assertion.
\end{proof}

\begin{remark}[The algebraic-factor refinement and its subsequent resolution]
\label{rem:DE-A3-mixed-fixed-refinement}
The specialization argument alone does not infer the formula for every
ordinary elliptic curve over $\overline{\mathbf F}_2$.
Theorem~\ref{thm:DE-A3-ppav-ranks} subsequently proves the formula
for all these factors, in every degree, as part of its universal
principal-surface result.  For the explicit algebraic
factor $y^2+xy=x^3+1$, exact retained Cartier calculations attain
the predicted bounds in degrees $1,\ldots,6$, and in every
rectangular system $(a,3)$ with $3\le a\le31$; the retained run
took $0.7093895580001117$ seconds.  The then-outstanding uniform
construction and failed sparse-section route remain recorded in
\texttt{notes/A3\_mixed\_rank\_research.md}.  These finite probes
are not substituted for the specialization proof or for the
later uniform proof.
The separate auxiliary certificate
\path{agents/obstructions_crystals/code/A3_mixed_generic_checks.py}
checks the integral Weierstrass invariants, their characteristic-two
reductions, the nonvanishing denominator in the $j$-invariant argument,
and the exact genus/rank identities for $10{,}000$ degrees.  It passed
in $0.04470579700137023$ seconds; its complete output is retained
beside the program.  These numerical identities support the displayed
construction, whereas the proof above supplies its uniform geometric
existence assertion.
\end{remark}
